%%%%%%%%%%%%%%%%%%%%%%%%%%%%%%%%%%%%%%%%%%%%%%%%%%%%%%%%%%%%%%%%%%%%%%
\section{Results}\label{sec:results}
%%%%%%%%%%%%%%%%%%%%%%%%%%%%%%%%%%%%%%%%%%%%%%%%%%%%%%%%%%%%%%%%%%%%%%

\subsection{Two decades of footprint growth and a hand-off between consumers}
\label{sec:res-destinations}

The land-use footprint of Brazilian soy embodied in consumption grew from
9.4~Mha in 2001 to 41.0~Mha in 2022, closely tracking the expansion of the
harvested area itself (Fig.~\ref{fig:destination}a). The defining structural
change of the period is a hand-off between the two dominant consumers
(Fig.~\ref{fig:destination}b). In 2001 the EU-27 accounted for 54\% of the
footprint and China for 11\%; the two curves cross in 2009--2010, and by 2022
China alone accounts for 42\% (17.2~Mha) against an EU-27 share below 20\%
(7.8~Mha). The EU footprint is remarkably stable in absolute terms
(5.1~Mha in 2001, 7.8~Mha in 2022): virtually all of the growth in Brazilian
soy land over two decades served demand growth in China and the rest of Asia
(1.3 to 12.0~Mha), not in Europe.

\begin{figure}[htbp]
\centering
\includegraphics[width=\textwidth]{fig_by_destination.pdf}
\caption{Brazilian soy land-use footprint by destination region, 2001--2022.
(a)~Absolute footprint (million hectares of soy land embodied in consumption).
(b)~Regional shares of the total footprint. Years before 2010 are computed on
the FABIO~v1.1 backend and 2010 onwards on FABIO~v2 (Methods); the vintage
seam shifts the Brazil-domestic share by a few percentage points around
2009/2010 but does not affect the ranking or the crossing point.}
\label{fig:destination}
\end{figure}

\subsection{From commodity exports to embodied animal products}
\label{sec:res-enduse}

Throughout the period, most of the footprint reaches consumers in commodity
form, as beans, soybean meal, or oil (Fig.~\ref{fig:enduse}). From 2010
onwards, when the end-use attribution is fully comparable across years,
roughly a quarter of the footprint is embodied in animal products
(meat, dairy and eggs) consumed at the destination, and the non-food and
industrial share roughly doubles after Brazil's 2008 biodiesel blending
mandate: the non-food footprint rises from 0.8~Mha in 2007 to 1.5~Mha in
2008 and 2.5~Mha by 2013. Shares before 2010 are shown for completeness but
are not directly comparable, because the v1.1 balance sheets book most
soy-chain consumption at the commodity level (Methods).

\begin{figure}[htbp]
\centering
\includegraphics[width=\textwidth]{fig_enduse.pdf}
\caption{Brazilian soy land-use footprint by end-use, 2001--2022.
(a)~Absolute footprint. (b)~End-use shares. The step at 2009/2010 in the
animal-products share is a classification seam between the FABIO v1.1 and v2
balance sheets, not a real dietary shift (Methods); end-use statements in the
text are therefore restricted to 2010--2022.}
\label{fig:enduse}
\end{figure}

\subsection{The footprint marches northwest}
\label{sec:res-spatial}

Spatially, the two decades redraw the map of where consumed soy land sits
(Fig.~\ref{fig:maps}). In 2001 the footprint concentrated in the southern
states and an emerging Mato Grosso core. By 2022 the centre of mass has moved
decisively into the Cerrado frontier: northern Mato Grosso and the MATOPIBA
region dominate, and the municipal pattern of China's footprint, nearly
empty in 2001 (1.0~Mha), spans the entire frontier by 2022 (17.2~Mha). The
EU-27 pattern, by contrast, changes far less, consistent with its flat
absolute footprint.

\begin{figure}[htbp]
\centering
\includegraphics[width=0.78\textwidth]{map_2001_2022.png}
\caption{Municipal origin of the Brazilian soy land-use footprint, 2001
versus 2022, for the total and the four major destination groups. All panels
share one colour scale (square-root scaled).}
\label{fig:maps}
\end{figure}

\subsection{Validation: agreement with Trase grows with flow size}
\label{sec:res-validation}

We validate the modelled municipality-to-country export flows against the
independent, shipment-based Trase supply-chain data (2004--2022).
Pooling all 127{,}862 municipality~$\times$~destination~$\times$~year flows
above one tonne, the log-flow correlation with Trase is $r = 0.68$ and 54\%
of flows fall within a factor of three of the Trase value. Agreement rises
strongly with flow size (Fig.~\ref{fig:rocket}): small flows scatter widely,
while the large flows that carry the footprint converge onto the 1:1 line,
so the model is most reliable exactly where the volume is. This holds in
every year, including the pre-2010 extension: the 2005 cross-section
validates at $r = 0.65$, statistically indistinguishable from 2010
($r = 0.73$). At the destination level, the pooled direct Pearson
correlation with Trase rises from 0.43 in 2004 to around 0.7 from 2011
onwards, and our 2013 value (0.685) reproduces the published benchmark of
the original single-year model (0.692). Per-tonne land intensities are
essentially flat across destinations (about 0.33~ha per tonne of
soybean-equivalent, the inverse of the national yield), with the EU-27 and
China within $\pm$4\% of each other over 2010--2015.

\begin{figure}[htbp]
\centering
\includegraphics[width=0.72\textwidth]{rocket_plot.pdf}
\caption{Rocket plot: modelled municipal soy export flows against Trase,
all years pooled (2004--2022), both axes log10 tonnes. Contours show point
density; the dashed line is 1:1. Flows below one tonne on either side are
excluded; the visible lower bound on the Trase axis (about 28~t) is the
granularity floor of the Trase composite itself.}
\label{fig:rocket}
\end{figure}
